\chapter{El Secreto del Director: Proyecto IZAR}

\begin{CajaSecreto}
\lead{CONFIDENCIAL // NIVEL DIRECTIVO OREKAK-CORP}\par
Este capítulo contiene la verdad oculta tras el Programa Azul y la conexión directa con el futuro distópico de NeoBilbao 2102. Los jugadores no deben tener acceso a esta información hasta que la descubran por sus propios medios durante la investigación.
\end{CajaSecreto}

\section{La farsa corporativa y la arquitectura del engaño}

\neofigure{03_proyecto_izar_laboratorio.jpg}{El búnker subterráneo del Proyecto IZAR: tanques amnióticos y cables de fibra óptica.}

Para la población civil de La Margen, el Programa Azul es el único faro de esperanza social. Orekak-Corp gasta millones de créditos en propaganda, concursos académicos y becas técnicas para convencer a la clase media de que el mérito personal puede comprar un billete hacia la Cúpula Dorada.

La realidad es una carnicería biotecnológica:

\begin{itemize}
  \item \lead{La verdadera función del Índice Azul:} Las pruebas psicotécnicas, los análisis de sangre en centros de salud y los tests de reflejos no miden la inteligencia ni la lealtad cívica. Evalúan parámetros biomédicos muy específicos: latencia de conducción sináptica, densidad proteica del líquido cefalorraquídeo y tolerancia al estrés neurotóxico.
  \item \lead{El objetivo de Orekak-Corp:} En 2084, los microprocesadores de silicio han tocado techo físico y los superconductores de Burdin-Urdin son aún inestables a temperatura ambiente. La corporación necesita una potencia de cálculo masiva para gestionar las redes de energía de la colonia y predecir los movimientos sociales. El Proyecto IZAR busca construir el primer superordenador biológico híbrido (\textbf{ADUR-OS v1.0}), utilizando cerebros humanos vivos como unidades de procesamiento en paralelo.
  \item \lead{El destino de los ganadores:} Los «premiados» no reciben un ático en la Cúpula Dorada. Son conducidos a un complejo blindado en el subsuelo de la Torre Orekak, sedados, rapados e intervenidos mediante craneotomías invasivas para injertarles puertos de fibra de cuarzo y cobre en el tronco encefálico.
\end{itemize}

\section{El defecto fatal: La Resonancia Azul}

El proyecto está condenado al colapso por una barrera biológica insalvable que el Director de Investigación, el \textbf{Dr. Victor Kessler}, se niega a admitir:

\begin{quote}
\textit{«El cerebro adulto posee una arquitectura identitaria rígida. Tiene recuerdos biográficos, traumas, filias, rencores y un ego impermeable. Cuando forzamos a dieciséis mentes adultas a sincronizarse en un mismo bus de datos, no cooperan: se devoran entre sí.»}
\flushright --- Dr. Bittor Albéniz, investigador desertor del Proyecto IZAR.
\end{quote}

Cuando la red neuronal supera el 75\% de carga de procesamiento, se desata el \textbf{Evento de Resonancia Azul}:
\begin{enumerate}
  \item \lead{Contaminación mnemónica:} Los sujetos empiezan a experimentar los recuerdos de los demás como si fueran propios (un oficinista recuerda el parto de otra mujer; una ingeniera canta las canciones de cuna del abuelo de un celador).
  \item \lead{Psicosis sináptica:} Las señales de rechazo provocan convulsiones eléctricas incontroladas que viajan a través del cableado del servidor, generando campos de estática que funden los monitores CRT y las lámparas halógenas.
  \item \lead{Fritura neuronal:} Si la conexión no se interrumpe físicamente, el fluido de los tanques hierve y las cortezas cerebrales sufren necrosis térmica instantánea.
\end{enumerate}

\section{La huida de Ane Garate}

Ane Garate fue integrada en el \emph{Cluster Theta} hace seis semanas. Su altísima resistencia sináptica le permitió mantener un hilo de consciencia individual mientras sus compañeros de nodo colapsaban. 

Durante una purga de emergencia causada por un cortocircuito, Ane fue desconectada y arrojada al vertedero de sujetos inviables por un celador sobornado o compasivo. Desorientada, cruzó las tuberías de desagüe de la Ría hasta emerger en su antiguo barrio. 

Su fuga es una amenaza existencial para Orekak-Corp: si se descubre que los «ganadores» son despedazados en tanques clandestinos, La Margen estallará en huelgas generales incontrolables.

\section{La conexión histórica con 2102: La semilla del Velo}

\begin{CajaRegla}{Revelación para la continuidad del universo}
La intervención de los personajes en 2084 cambiará el curso de la historia de NeoBilbao:
\begin{itemize}
  \item Al destruir o sabotear el Proyecto IZAR en la Fase 4, los jugadores fuerzan a la junta directiva de Orekak-Corp a archivar el uso de cerebros adultos.
  \item La conclusión de los científicos corporativos será aterradora: \emph{«El problema no es la red biológica; el problema es la edad. Los adultos tienen demasiado pasado. Necesitamos mentes dúctiles, moldeables y puras. Mentes infantiles»}.
  \item Esta catástrofe de 2084 es la chispa directa que, dieciocho años más tarde (2102), impulsará a Orekak-Corp a crear el \textbf{Etxearen Zina} (el Juramento del Linaje) y el \textbf{Velo}: arrebatar a los hijos de los mineros para utilizarlos como nodos del superordenador Mari.
\end{itemize}
\end{CajaRegla}
